\scanpage{86}
% archive foliation: 85

On peut préciser, en introduisant l'algèbre engendrée par
\deleted{\unsure{$A$}} $G(\mathbb Q)$ dans $A$, soit $B$.
On sait que $B=\prod L_i$, chaque $L_i$ ext. quadratique d'un corps
de nb $K_i$ tt. réel.
($K_i$ déterminé \unsure{uniquement}
\deleted{\unsure{par la condition}}\inserted{dans $L_i$}
comme \unsure{l'un des s/c tt. réels} de $L_i$),
l'involution est donnée sur $L_i$ par la conjugaison p.r. à $K_i$
(i.e. la conjugaison complexe ordinaire, si $L_i\subset\mathbb C$),
la trace est proportionnelle, sur chaque $L_i$, à la trace ordinaire.
Ainsi $B^*$ n'est autre que $H(\mathbb Q)$, où
\[
H=\prod_i\left(\prod_{L_i/\mathbb Q}G_{m\,L_i}\right)
\]
et le \unsure{groupe|gr.} unitaire $B^{*u}$ n'est autre que \unsure{le} \deleted{\unsure{groupe}} des pts rationnels de
\[
H^u=\prod_i\left(\prod_{K_i/\mathbb Q}\unsure{\Gamma_i|\mathcal T_i}\right)
\]
où $\unsure{\Gamma_i|\mathcal T_i}$ est la forme tordue de $G_m$ sur $K_i$
définie par l'extension quadratique $L_i$ de $K_i$.
\deleted{[}\deleted{Ainsi}
Cela met en évidence que $H^u$ est un tore ``compact'',
\deleted{\unsure{donc il est}}
(de nature d'ailleurs assez particulière!), et aussi $U^\circ$ comme
sous-tore de ce dernier \ldots

NB. Il convient de déterminer $\varepsilon:G\longrightarrow G_m$
en termes des \unsure{hom. normes} sur les $L_i$, et
$j:G_m\longrightarrow G$ \inserted{de même \ldots}
